\documentclass[11pt,a4paper]{ctexart}
\usepackage[margin=2.6cm]{geometry}
\usepackage{amsmath,amssymb,bm}
\usepackage{booktabs}
\usepackage{enumitem}
\usepackage{setspace}
\usepackage{xcolor}
\usepackage{hyperref}
\usepackage{tcolorbox}
\tcbuselibrary{breakable,skins}
\usepackage{fancyhdr}
\usepackage{titlesec}
\usepackage{microtype}

\hypersetup{
  colorlinks=true,
  linkcolor=blue!50!black,
  urlcolor=blue!60!black,
  citecolor=blue!50!black
}

\definecolor{lockbg}{HTML}{E8F5E9}
\definecolor{lockfr}{HTML}{2E7D32}
\definecolor{openbg}{HTML}{FFF8E1}
\definecolor{openfr}{HTML}{F9A825}
\definecolor{laterbg}{HTML}{ECEFF1}
\definecolor{laterfr}{HTML}{546E7A}
\definecolor{storybg}{HTML}{E3F2FD}
\definecolor{storyfr}{HTML}{1565C0}

\newtcolorbox{lockbox}[1][]{
  colback=lockbg,colframe=lockfr,title={已锁定},
  fonttitle=\bfseries,breakable,sharp corners=southwest,#1
}
\newtcolorbox{openbox}[1][]{
  colback=openbg,colframe=openfr,title={还没锁，下一轮一起拍},
  fonttitle=\bfseries,breakable,sharp corners=southwest,#1
}
\newtcolorbox{laterbox}[1][]{
  colback=laterbg,colframe=laterfr,title={以后再做，现在故意不做},
  fonttitle=\bfseries,breakable,sharp corners=southwest,#1
}
\newtcolorbox{storybox}[1][]{
  colback=storybg,colframe=storyfr,title={这一节在讲什么},
  fonttitle=\bfseries,breakable,sharp corners=southwest,#1
}

\newcommand{\TabUF}{\textsf{TabUF}}
\newcommand{\Disco}{\textsf{DiscoSCM}}
\newcommand{\uu}{\mathbf{u}}
\newcommand{\UU}{\mathbf{U}}
\newcommand{\WW}{\mathbf{W}}
\newcommand{\EE}{\mathbf{E}}
\newcommand{\YY}{\mathbf{Y}}

\setlength{\headheight}{14pt}
\pagestyle{fancy}
\fancyhf{}
\lhead{\TabUF{} 数据生成活文档}
\rhead{v0.1 \quad 2026-08-14}
\cfoot{\thepage}
\renewcommand{\headrulewidth}{0.3pt}

\setstretch{1.18}
\setlist[itemize]{leftmargin=1.4em,itemsep=0.25em,topsep=0.3em}
\setlist[enumerate]{leftmargin=1.6em,itemsep=0.25em,topsep=0.3em}

\title{\TabUF{} 数据生成技术文档\\[0.4em]
\large 从总体到 Episode：一篇会跟着我们一起长的说明书}
\author{Heyang Gong \and 文献调研（协作草稿）}
\date{2026 年 8 月 14 日 \quad 草稿 v0.1}

\begin{document}
\maketitle
\vspace{-1.2em}
\begin{center}
\textit{这不是最终 API 规格。API 会跟着这篇文档一起改。}\\[0.3em]
{\small 仓库：\url{https://github.com/1587causalai/tabuf-episode-api}}
\end{center}

\tableofcontents
\newpage

\section{先把话说清楚}

\begin{storybox}
我们现在最缺的不是又一个网络结构，而是\textbf{可训练的 episode 从哪来}。
这篇文档只做一件事：把「数据怎么生成」讲到你可以抬杠、可以改、可以写成代码的程度。
HTTP 路径、JSON 字段名、端口号，都是后话。
\end{storybox}

你已经有一个能跑的 \texttt{POST /v0/episodes}。它远远不够，原因很简单：
它把第零版生成律塞进了一个临时接口里，但没有把\textbf{生成这件事本身}讲清楚。
接口可以换，生成语义不能糊。

所以我们换一个工作方式。以后每一轮讨论，都先改这篇文档，再改代码。
文档里用三种颜色盒子：

\begin{itemize}
  \item \textbf{绿：已锁定。} 改它等于改项目方向，要明确说「我反悔了」。
  \item \textbf{黄：还没锁。} 这一轮可以拍，也可以先挂着。
  \item \textbf{灰：以后再做。} 不是忘了，是第零版故意关掉，免得引擎一开始就长歪。
\end{itemize}

\begin{lockbox}
第零版只生成\textbf{观测世界}的 Unit$\times$Feature 格子，再切成格子级的上下文 $C$ 与查询 $Q$。
调用方拿到的是 episode，不是 DAG，也不是潜变量 $\uu_i$。
\end{lockbox}

\section{我们到底要喂给网络什么}

表格基础模型的训练对象，常常被说成「一张表」。
这对 \TabUF{} 不够用。一张完整的表，更像一次人口普查的结果；
网络真正要学的，是：\textbf{在已经看见一部分格子之后，去填另一部分格子}。

所以训练样本不是表，是 \textbf{episode}：

\begin{equation}
  \text{episode}
  \;=\;
  \bigl(\,Y^{\mathrm{in}},\; M^{C},\; M^{Q},\; y^{Q}\,\bigr).
\end{equation}

记号先定下来，后面一直用：

\begin{itemize}
  \item $n$：这一次实现里有多少个 unit（行）。
  \item $d$：有多少个 feature（列）。
  \item 格子 $(i,j)$：第 $i$ 个 unit 在第 $j$ 个 feature 上的取值 $Y_{ij}$。
  \item $C$：上下文格子。网络可以看见这些值。
  \item $Q$：查询格子。网络要预测这些值。第零版只训回归臂 $R$。
  \item $Y^{\mathrm{in}}$：把 $Q$ 位置挖空之后的输入表。
  \item $y^{Q}$：按 $M^{Q}$ 拉成向量的查询真值。
\end{itemize}

表 $\YY$ 是完整的，不因 mask 挖空。
$M$ 与 $Q$ 独立按全表比例抽取，允许重合：
\begin{equation}
  Q\cap M \text{ 缺失填补},\qquad
  Q\setminus M \text{ 普通预测}.
\end{equation}
没有单独的 $y$ 标签。上下文是 $\neg M\wedge\neg Q$。
默认 $P(M)=0.05$、$P(Q)=0.15$，并集不必是 $20\%$。
每条 episode 各自抽一个总体。
列类型权重由 API 的 \texttt{type\_weights} 传入，默认 $70{:}5{:}10{:}5{:}5$（数值/有序/二值/多值/超多值）；\texttt{independent\_frac} 默认 $5\%$。

为什么是格子级，而不是整行、整列当 query？
因为 \TabUF{} 要能回答的，不只是「这个人缺一个标签」，
还包括「这个特征在这些人身上会是什么」。
推荐场景里，这几乎是默认形态：人 $\times$ 物品 的一部分被看见，另一部分要估。

\begin{lockbox}
训练样本是格子级 $C/Q$ 的 episode，不是整张表，也不是一个 DAG。
$C\cap Q=\emptyset$，第零版 $C\cup Q$ 覆盖全表。
\end{lockbox}

\section{为什么必须先抽总体，再填格子}

这一节是整篇文档的脊柱。如果你只记住一件事，记住它。

\subsection{行不是一个新世界}

常见的表格合成器（包括不少 SCM / 树 SCM 预训练引擎）骨子里是：
\textbf{每一行独立地再抽一遍机制}。
行与行之间可以相关，但「这一行是谁」往往没有单独的身份。

\Disco{} 不这么看。它把随机性劈成两截 \cite{gong2024discoscm}：

\begin{itemize}
  \item $U$：unit selection。一次实现 $U=u$ 就是一个个体。
  \item $\mathbf{E}$：这个个体身上的外生噪声。它独立于 $U$。
\end{itemize}

结构方程是 \textbf{unit-specific} 的：

\begin{equation}
  v_{\ell} \;\leftarrow\; f_{\ell}(pa_{\ell},\, e_{\ell};\, u).
  \label{eq:unit-law}
\end{equation}

同一个机制 $f_{\ell}$，换一个 $u$，响应可以不同。
$u$ 钉住的是个体，不是行号，更不是「又一张新表」。

于是生成顺序被钉死了：

\begin{enumerate}
  \item 先有一个总体 $\Pi$，从中实现 $n$ 个 unit，得到 $\{\uu_i\}_{i=1}^{n}$。
  \item 再对这些已经实现的人，按 \eqref{eq:unit-law} 填他们在各个 feature 上的取值。
\end{enumerate}

\texttt{sample\_population} 然后 \texttt{fill\_grid}。
不能倒过来，也不能把「抽一行噪声」误写成「抽一个新总体」。

\subsection{一个小例子，把这句话听进去}

假设总体里有两种人。第一种人对折扣很敏感，第二种人几乎不买账。
如果你每一行都重新抽一套机制，表会看起来很热闹，但网络学到的是
「行与行之间的噪声花样」，不是「同一种人在不同列上应当怎么响应」。

\TabUF{} 要的是后者。
推荐、Uplift、个性化决策，问的都是：\textbf{这个 unit 换一个 feature / 换一个处理，会怎样}。
没有钉住的 $u$，后面的反事实臂根本没有主语。

\begin{lockbox}
先从总体 $\Pi$ 实现 $\{\uu_i\}$，再用 unit-specific response law 填格子。
$\uu_i$ 是个体，不是 row-id。不要把每一张表、每一行当成一个新总体。
\end{lockbox}

\begin{laterbox}
第零版\textbf{不}做干预，\textbf{不}重抽反事实噪声 $\mathbf{E}(\mathbf{x})$。
那是 \Disco{} Layer~3 的能力，也是以后 S 臂 / 反事实 episode 的事。
现在只抽一次观测噪声 $\EE$，得到一个观测世界。
\end{laterbox}

\section{三层流水线：Prior / World / Compiler}

把生成想成三条传送带，比想成「一个函数吐 JSON」清楚得多。
网络不准上传送带。传送带只负责世界和考题。

\subsection{Prior：抽总体，不抽考题}

Prior 决定「有哪些人、列的方向是什么」。
第零版 Prior 非常瘦：

\begin{align}
  \uu_i &\sim \mathcal{N}(0, I_k),
  && i=1,\ldots,n,
  \label{eq:prior-u}\\
  \mathbf{w}_j &\sim \mathcal{N}(0, I_k),
  && j=1,\ldots,d,
  \label{eq:prior-w}
\end{align}

然后（可选）把 $\WW\in\mathbb{R}^{d\times k}$ 的每一列做 $\ell_2$ 归一，
让 $k$ 个潜在轴在特征空间里尺度差不多。
$\{\uu_i\}$ 合起来就是这一次实现的总体切片；
$\WW$ 是这个世界上共享的 feature 方向。

这里的 $k$ 叫 \texttt{unit\_dim}。它不是「网络宽度」，是个体嵌入的生成维数。

\subsection{World：按响应律把格子填满}

观测世界第零版只承认一条线性因子律：

\begin{equation}
  Y_{ij}
  \;=\;
  \langle \uu_i,\, \mathbf{w}_j \rangle
  + e_{ij},
  \qquad
  e_{ij}\sim\mathcal{N}(0,\sigma^2).
  \label{eq:factor}
\end{equation}

矩阵写法更干净：

\begin{equation}
  \YY \;=\; \UU\,\WW^{\top} + \EE,
  \qquad
  \UU\in\mathbb{R}^{n\times k},\;
  \WW\in\mathbb{R}^{d\times k},\;
  \EE\in\mathbb{R}^{n\times d}.
\end{equation}

这不是 MiniSCM，也不是 TabICL 的 ExtraTrees 生成器，
更不是把 TabPFN 未公开先验包一层皮。
它是 \Disco{} 语义在表格格子上的\textbf{最瘦实现}：
个体 $\uu_i$ 与列方向 $\mathbf{w}_j$ 的内积，加上一次观测噪声。

为什么从这么瘦的东西开始？
因为我们要先验证三件事能不能同时成立：

\begin{enumerate}
  \item 同一个人在不同列上的取值，真的被同一个 $\uu_i$ 串起来；
  \item 同一列在不同人身上的取值，真的被同一个 $\mathbf{w}_j$ 串起来；
  \item 网络只看见格子，却仍能把查询格填对。
\end{enumerate}

这三件事过不了，后面换更漂亮的 $f_j(\uu_i)$ 没有意义。

\subsection{Compiler：出题，不改世界}

World 给出完整的 $\YY$。Compiler 不碰生成律，只出题：

\begin{enumerate}
  \item 在 $n\cdot d$ 个格子里随机抽约 $q$ 比例进 $Q$，其余进 $C$。
  \item 做 $Y^{\mathrm{in}}$：把 $Q$ 位置写成「看不见」（实现上是 \texttt{NaN} / JSON \texttt{null}）。
  \item 把 $Q$ 上的真值按行优先拉成 $y^{Q}$。
  \item \textbf{默认丢掉} $\UU$ 和 $\WW$。
\end{enumerate}

Compiler 还可以顺手带一个 S 臂诊断位 $\alpha$（例如「这一格是否该被当成反事实查询」）。
第零版\textbf{可以发 $\alpha$，但不准拿它和 $R$ 损失平均成一个头}。
现在只训 $R$。

\begin{lockbox}
流水线固定为 Prior $\to$ World $\to$ Compiler。
网络不进生成器。$R$ 与 $S$ 即使将来同 episode，也不许静默平均成一个头。
真目标表上从第一天就冻 $\theta$，合成引擎只负责预训练 episode。
\end{lockbox}

\begin{openbox}
下列细节还没锁，正好是后面几轮要一起拍的：
\begin{itemize}
  \item 同一个总体能否跨多个 episode 复用（先抽一个大 $N$ 的池子，再每次取 $n$ 行）？
  \item $\WW$ 是每个 episode 独立抽，还是一个「世界」共享、episode 只换 query 掩码？
  \item $q$ 固定 $0.15$，还是对 $q$ 本身做 curriculum？
  \item 列归一化是默认开，还是变成 Prior 的一个开关？
\end{itemize}
\end{openbox}

\section{第零版生成律，逐步拆开}

把 \eqref{eq:factor} 再讲一遍，这次按代码会执行的顺序。

\paragraph{第一步：实现总体。}
抽 $\UU\in\mathbb{R}^{n\times k}$。
行 $\uu_i$ 就是第 $i$ 个人。
不要在注释里写 ``sample a table of rows''。写 ``sample a population of units''。

\paragraph{第二步：实现列方向。}
抽 $\WW\in\mathbb{R}^{d\times k}$。
行 $\mathbf{w}_j$ 是第 $j$ 列在个体空间里的方向。
列归一化（如果开）是

\begin{equation}
  \mathbf{w}_{:r}
  \;\leftarrow\;
  \mathbf{w}_{:r}
  \big/
  \lVert \mathbf{w}_{:r} \rVert_2,
  \qquad r=1,\ldots,k.
\end{equation}

\paragraph{第三步：一次观测噪声。}
$\EE\sim\mathcal{N}(0,\sigma^2 I)$，与 $\UU$ 独立。
这就是 \Disco{} 里的 factual noise。第零版到此为止，不再抽 $\EE(\mathbf{x})$。

\paragraph{第四步：填满格子。}
$Y_{ij}=\langle \uu_i,\mathbf{w}_j\rangle + e_{ij}$。
此刻世界已经完成。后面任何掩码都不许回头改 $Y_{ij}$。

\paragraph{第五步：出题。}
令 $N_{\mathrm{cell}}=nd$，
$n_Q=\mathrm{clip}\bigl(\mathrm{round}(q\cdot N_{\mathrm{cell}}),\,1,\,N_{\mathrm{cell}}-1\bigr)$，
均匀无放回抽 $Q$。
$C$ 是补集。于是一定有至少一个上下文格、至少一个查询格。

\paragraph{第六步：打包给调用方。}
默认 payload 只有世界的可观察部分和考题，没有潜变量。

\begin{lockbox}
第零版响应律就是 \eqref{eq:factor}。
数值、加性、高斯噪声、观测一次。
不把 TabICL / TabPFN 先验混进同一套 epoch。
\end{lockbox}

\section{调用方能看见什么，不能看见什么}

这一节是网络合同。API 字段名以后可以改，合同不能偷偷改。

默认返回：

\begin{center}
\begin{tabular}{@{}llp{8.2cm}@{}}
\toprule
字段 & 形状 & 含义 \\
\midrule
\texttt{y\_input} & $n\times d$ & 上下文格是 $Y_{ij}$，查询格是空 \\
\texttt{context\_mask} & $n\times d$ & $M^{C}_{ij}=1$ 当且仅当 $(i,j)\in C$ \\
\texttt{query\_mask} & $n\times d$ & $M^{Q}_{ij}=1$ 当且仅当 $(i,j)\in Q$ \\
\texttt{y\_query} & $n_Q$ & $Y$ 在 $Q$ 上按行优先拉直 \\
\texttt{shapes} & --- & 上面这些张量的形状，避免客户端猜 \\
\texttt{seed} & 标量 & 这一集的名字，可复现 \\
\bottomrule
\end{tabular}
\end{center}

默认\textbf{不}返回 $\UU$、$\WW$、完整 $\YY$。
\texttt{debug=true} 才附带，而且 debug 输出不准进训练 DataLoader。

为什么这么凶？
因为 $\uu_i$ 一旦进网络，模型会走捷径：直接读个体真名，不再从格子里推断「这个人是谁」。
预训练就会变成漏题。真表上你也没有 $\uu_i$ 可喂。

\begin{lockbox}
训练路径禁止 $u_i$ 进网。
Compiler 的张量合同先于 Prior 插件化。
小格子 JSON，大格子以后再换 \texttt{npz}，但语义不变。
\end{lockbox}

\begin{openbox}
张量合同里还没锁的：
\begin{itemize}
  \item 要不要带列类型 \texttt{column\_types}（第零版全是实数，字段可以先占位）？
  \item $y^{Q}$ 的拉直顺序：行优先已经写进代码，是否正式写进合同？
  \item 多 episode 一次返回时，是列表还是 padded batch？
  \item S 臂诊断 $\alpha$ 的形状：格子级，还是 episode 级？
\end{itemize}
\end{openbox}

\section{当前 API 草稿（会变）}

把现在仓库里的接口抄在这里，只作为\textbf{实现快照}，不是承诺。

\texttt{GET /health} 返回版本号。

\texttt{POST /v0/episodes} 现在接受：

\begin{center}
\begin{tabular}{@{}lcp{7.6cm}@{}}
\toprule
参数 & 第零版默认 & 备注 \\
\midrule
\texttt{n\_units} & 64 & 这一集实现多少 unit \\
\texttt{n\_features} & 8 & 多少列 \\
\texttt{unit\_dim} & 4 & $k$，个体潜空间维 \\
\texttt{query\_frac} & 0.15 & 要预测的格子，与 missing 可重合 \\
\texttt{missing\_frac} & 0.05 & 世界缺失，与 query 可重合 \\
\texttt{sigma} & 0.3 & 观测噪声标准差 \\
\texttt{seed} & 0 & 第 $e$ 集用 $\mathrm{seed}+e$ \\
\texttt{n\_episodes} & 1 & 单次最多 32，纯属保护网关 \\
\texttt{return\_mechanism} & false & 训练必须 false；true 则返回 DiscoSCM \\
\texttt{debug} & false & 训练必须 false；隐含上一开关并附满表 \\
\bottomrule
\end{tabular}
\end{center}

这些数字没有神圣性。$64\times 8$ 只是为了能用手印 JSON。
真正训练时会走 $n,d$ curriculum：先小后大，再 pad 到本 epoch 的 $\max n\times \max d$。

\begin{openbox}
API 形态本身就是开放问题，也是这份文档存在的理由。例如：
\begin{itemize}
  \item 继续 REST JSON，还是一开始就按 DataLoader 协议走共享内存 / \texttt{npz}？
  \item 要不要 \texttt{population\_id}，让多次请求钉在同一次总体实现上？
  \item 版本号是 URL 里的 \texttt{/v0}，还是 payload 里的 \texttt{schema} 字段？
  \item 失败时如何区分「参数不合法」和「生成器内部崩了」？
\end{itemize}
我们不假装现在就能选对。先把生成语义聊稳，接口跟着语义长。
\end{openbox}

\section{刻意关掉的东西}

写下来，免得下一周手痒又加回来。

\begin{laterbox}
\begin{itemize}
  \item 干预 $do(\mathbf{x})$，以及 counterfactual noise 重抽。
  \item 自然缺失（World 级 missingness）。现在的空格只来自 Compiler 出题。
  \item 类别列、混合类型、文本单元格。
  \item S 臂训练目标。诊断 $\alpha$ 可以留位，损失不能混。
  \item 把 TabICL ExtraTrees / 未公开 TabPFN 先验接到同一套 epoch。
  \item 可插拔节点函数 $f_j$（MLP、树、周期、跳跃）。Prior 插件化要等张量合同先冻住。
  \item 把网络、tokenizer、损失写进这个仓库。这是数据引擎，不是训练引擎。
\end{itemize}
\end{laterbox}

关掉不是因为它们不重要。
恰恰相反：它们每一个都会逼我们改 World 的语义。
语义没稳之前加功能，等于在流沙上盖楼。

\section{一个完整的数值散步}

假设 $n=4$，$d=3$，$k=2$，$q=0.25$，$\sigma=0.3$，\texttt{seed}=0。
（数字是示意，不以某次运行的浮点为准。）

\begin{enumerate}
  \item Prior 抽出 4 个 unit，每个是 $\mathbb{R}^{2}$ 里的一个点。
        这 4 个人组成这一集的总体切片。
  \item 再抽出 3 个列方向，同样在 $\mathbb{R}^{2}$ 里。
        列 1 也许更对齐第一个潜在轴，列 2 更对齐第二个。
  \item 对每个人、每一列做内积，再加一点点高斯噪声，得到 $4\times 3$ 的满表。
  \item $12$ 个格子里大约 $3$ 个进 $Q$。网络看见另外 $9$ 个，去猜这 $3$ 个。
  \item 调用方拿到挖空的表和三个真值。它不知道这 4 个人的 $\uu_i$，
        也不知道三列的 $\mathbf{w}_j$。
\end{enumerate}

如果网络真的在学 \TabUF{} 该学的东西，它必须从那 $9$ 个可见格子里，
同时推断「人」和「列」，再在交点上给出预测。
这就是因子律作为第零版考题的原因：答案结构清楚，作弊路径也清楚
（把 $\UU$ 漏出去立刻满分）。

\section{以后会变复杂的方向（只指路，不开工）}

等第零版在合成格子上把 $R$ 臂打穿，再按这条路长，而不是按心情长：

\begin{enumerate}
  \item \textbf{World 级缺失。} 先有世界掉格子，再在剩下的格子上出题。
        $C\cup Q$ 不再等于全表。
  \item \textbf{非线性 unit-specific 律。}
        $Y_{ij}=f_j(\uu_i)+e_{ij}$，
        $f_j$ 从线性换成 MLP / 周期 / 分段，但仍是同一个 $u$ 贯穿一行。
  \item \textbf{干预 episode。}
        对某一列做 $do$，重抽 $\EE(\mathbf{x})$，Compiler 问反事实格。
        那时 S 臂才有真实的监督。
  \item \textbf{类别与混合类型。} 先锁实数合同，再扩。
  \item \textbf{$n,d$ curriculum 与无限 DataLoader。}
        生成器永远在线，pad 到本 batch 的最大格子。
\end{enumerate}

每加一层，先改这篇文档的对应小节，再改代码。
代码不许跑到文档前面去。

\section{下一轮我们要一起拍的问题}

把开放问题收成一张清单。你可以从里面点名先聊哪一条。

\begin{enumerate}[leftmargin=1.8em]
  \item 一个 episode 是否必须对应一次全新的总体实现？
        还是允许 \texttt{population\_id} 钉住同一批 $\uu_i$，只换掩码？
  \item $\WW$ 的生命周期：每集一抽，还是一个世界多次出题？
  \item $q=0.15$ 是暂时的手感，还是写成合同？
  \item 列类型字段现在占位，还是干脆第零版不出现？
  \item 无限 DataLoader 是这个仓库的事，还是训练仓库来拉 API？
  \item debug 通道要不要用单独的端口 / 路径，避免误进训练？
  \item 文档本身继续中文活文档，还是同时养一份英文（给以后写论文用）？
\end{enumerate}

我的建议（可以否决）：
第 1、2 题下一轮先拍，因为它们决定 Prior 要不要有状态；
第 5 题可以晚一点，等生成语义稳了再把运输层做厚。

\section*{这份文档怎么用}

\begin{itemize}
  \item 同意某一段：说「锁第 X 节」。我会把黄盒子改成绿盒子，再改代码。
  \item 反对某一段：直接改句子。不要只说「API 再想想」——指出是 Prior、World 还是 Compiler 错了。
  \item 想加功能：先在「以后再做」里找，确认它不是第零版故意关掉的。
        如果要提前做，写明为什么现在就必须做。
\end{itemize}

第零版代码只是这篇文档第 4--6 节的一个临时投影。
投影可以砸。语义要留下。

\begin{thebibliography}{9}
\bibitem{gong2024discoscm}
Heyang Gong, Chaochao Lu, Yu Zhang.
\newblock Distribution-consistency Structural Causal Models.
\newblock 2024.
\newblock 本地指南：\texttt{discoscm/latex/main-guide.tex}。

\bibitem{pearl2009}
Judea Pearl.
\newblock \emph{Causality}.
\newblock Cambridge University Press, 2nd edition, 2009.

\bibitem{holland1986}
Paul W. Holland.
\newblock Statistics and causal inference.
\newblock \emph{JASA}, 1986.
\end{thebibliography}

\end{document}
